\begin{mistake}{2026-07-21}{02}

\begin{problemcontent}
已知 \(y_1 = 3, y_2 = 3 + x^2, y_3 = 3 + e^x\) 是某二阶线性非齐次方程的三个特解，求该微分方程及通解。
\end{problemcontent}

\begin{solutioncontent}
由非齐次方程解的结构定理，非齐次特解之差为对应齐次方程的解。
可得齐次方程的两个线性无关解为：
\[ Y_1 = y_2 - y_1 = x^2, \quad Y_2 = y_3 - y_1 = e^x \]
取 \(y_1 = 3\) 为非齐次方程的一个特解，则原方程的通解为：
\[ y = C_1 x^2 + C_2 e^x + 3 \quad \text{①} \]

为求原微分方程，需通过求导消去任意常数 \(C_1, C_2\)：
\[ y' = 2C_1 x + C_2 e^x \quad \text{②} \]
\[ y'' = 2C_1 + C_2 e^x \quad \text{③} \]

利用 \(C_2 e^x\) 求导形式不变的特点，两两相减消去 \(C_2\)：
③式 \(-\) ②式 得：
\[ y'' - y' = 2C_1(1 - x) \quad \text{④} \]
①式 \(-\) ②式 得：
\[ y - y' = C_1(x^2 - 2x) + 3 \implies y - y' - 3 = C_1(x^2 - 2x) \quad \text{⑤} \]

利用交叉相乘消去 \(C_1\)：将④式乘以 \((x^2 - 2x)\)，⑤式乘以 \(2(1 - x)\)，两式相等：
\[ (x^2 - 2x)(y'' - y') = 2(1 - x)(y - y' - 3) \]

展开并移项合并 \(y, y', y''\) 的同类项，得到所求微分方程：
\[ (2x - x^2)y'' + (x^2 - 2)y' + 2(1 - x)y = 6(1 - x) \]
\end{solutioncontent}

\mistakeknowledge{
\begin{itemize}
  \item \textbf{核心定理}：非齐次线性微分方程的解的结构：非齐次特解之差为对应齐次方程的解。
  \item \textbf{消去常数法反推方程}：已知通解求微分方程，本质是利用连续求导及代数变形消去通解中的任意常数（如 \(C_1, C_2\)）。
  \item \textbf{解题技巧}：当通解中含有 \(e^x\) 时，优先利用“两式相减”直接消去带 \(e^x\) 的常数项，再通过交叉相乘法消去多项式常数项，可大幅减少计算量。
\end{itemize}}

\end{mistake}
